\chapter{Introduction}
\label{ch:intro}

\section{What pyCMF is}

\pycmf{} is a Python package of \emph{correlated mean-field methods} built on
\pyscf{}. Its core method is one-body second-order M\o ller--Plesset
perturbation theory (OBMP2). In OBMP2 the Hartree--Fock Fock operator is
replaced by a \emph{correlated Fock operator} that contains a one-body MP2
correlation potential. The orbitals, the orbital energies and the amplitudes
are then iterated to self-consistency, at the cost of one MP2 calculation per
cycle.

Because the correlated Fock operator is a one-body operator, the methods of
\pycmf{} behave like a mean-field theory in many respects. They give orbitals,
orbital energies and a single determinant, they can be started from a
non-Aufbau determinant to describe excited states, and they can be used to
dress the one-body part of an active-space Hamiltonian. On top of OBMP2 the
package provides:

\begin{itemize}
  \item O2BMP2, which keeps only the opposite-spin amplitudes and scales them;
  \item one-body double hybrids (OBDH, O2BDH), in which a Kohn--Sham Fock
        matrix is augmented by a scaled OBMP2 correlation potential;
  \item state-specific excited states with the maximum overlap method (MOM)
        and its initial variant (IMOM);
  \item one-body downfolding (OBDF), which produces an active-space
        Hamiltonian for any many-body solver;
  \item properties: densities, dipole moments, $\langle S^2\rangle$ and
        isotropic hyperfine coupling constants;
  \item experimental extensions: restricted open-shell OBMP2, projection
        embedding in DFT, and OBMP2 and downfolding for periodic systems with
        $k$-point sampling.
\end{itemize}

\section{Methods and their status}

Table~\ref{tab:methods} lists the methods of version~1.0. \emph{Stable}
methods are covered by the test suite and exported at the top level of the
package (\code{import pycmf}). \emph{Experimental} methods are under
development: they are imported from their sub-package, emit an
\code{ExperimentalWarning} when they are instantiated, and are described in
Chapter~\ref{ch:validation} together with their known issues.

\begin{table}[!htb]
\centering
\small
\caption{Methods of \pycmf{} 1.0.}
\label{tab:methods}
\begin{tabularx}{\textwidth}{@{}l X l l@{}}
\toprule
Method & Description & Entry point & Status \\
\midrule
OBMP2 & self-consistent one-body MP2, restricted and unrestricted & \code{pycmf.OBMP2} & stable \\
O2BMP2 & OBMP2 with scaled opposite-spin amplitudes only & \code{pycmf.O2BMP2} & stable \\
OBDH & one-body double hybrid: Kohn--Sham Fock matrix plus scaled OBMP2 potential & \code{pycmf.OBDH} & stable \\
O2BDH & OBDH with scaled opposite-spin amplitudes only & \code{pycmf.O2BDH} & stable \\
OBDF & one-body downfolding to an active-space Hamiltonian & \code{pycmf.OBDF} & stable \\
excited states & any of the first four with \code{occupation='mom'} or \code{'imom'} & keyword & stable \\
RO-OBMP2 & restricted open-shell OBMP2 & \code{pycmf.mol.roobmp2.ROOBMP2} & experimental \\
embedding & OBMP2- and OBDH-in-DFT projection embedding, concentric localization & \code{pycmf.embed} & experimental \\
KOBMP2, KO2BMP2 & periodic systems, closed shell: density fitting, $k$-point symmetry, MPI & \code{pycmf.pbc} & experimental \\
KOBDF & downfolding to active bands with $k$-points & \code{pycmf.pbc.KOBDF} & experimental \\
KUOBMP2 & periodic systems, unrestricted (code of version 0.2) & \code{pycmf.pbc.KUOBMP2} & experimental \\
\bottomrule
\end{tabularx}
\end{table}

\section{Conventions of this manual}

\begin{itemize}
  \item Energies are in Hartree unless stated otherwise. For periodic
        systems all energies are per unit cell.
  \item Indices $i,j,k,l$ denote occupied orbitals, $a,b,c,d$ virtual
        orbitals and $p,q,r,s$ general orbitals. Two-electron integrals are
        written in chemists' notation, $(pq|rs)$.
  \item Python code is shown in framed boxes. Lines that start with
        \code{$} are shell commands.
  \item Names in \texttt{typewriter} font are Python objects or files of the
        repository. Paths are relative to the root of the repository.
\end{itemize}

\section{Design principles}

\pycmf{} follows the conventions of \pyscf{}, so that a user of \pyscf{} can
use it without learning a new interface:

\begin{enumerate}
  \item A method object is built from a converged mean-field object of
        \pyscf{}. Options are attributes of the object.
  \item \code{kernel()} (or \code{run()}, which returns the object itself)
        performs the calculation and returns the main result, the total
        energy. All other results are attributes.
  \item The input mean-field object is never modified.
  \item Each method has exactly one implementation. Density fitting, DIIS,
        the occupation scheme, the spin scaling and the order of the
        Baker--Campbell--Hausdorff expansion are attributes of that
        implementation, not separate classes or modules.
\end{enumerate}

\section{How to cite}

If you use \pycmf{}, please cite the papers of the methods you use
(Appendix~\ref{ch:references}) and \pyscf{}. The file \code{CITATION.cff} in
the repository contains the same information in machine-readable form.

\begin{itemize}
  \item OBMP2: Lan and Yanai, J.~Chem.~Phys.~\textbf{138}, 224108 (2013);
        Tran, J.~Phys.~Chem.~A \textbf{125}, 9242 (2021).
  \item Densities: Tran, Phys.~Chem.~Chem.~Phys.~\textbf{24}, 19393 (2022).
  \item O2BMP2: Tran, Nguyen and Tran, J.~Phys.~Chem.~A \textbf{128}, 1543 (2024).
  \item State-specific excited states: Tran and Tran, J.~Chem.~Phys.~\textbf{162}, 104109 (2025).
  \item Downfolding: Le and Tran, J.~Phys.~Chem.~A \textbf{127}, 5222 (2023).
  \item \pyscf{}: Sun \emph{et al.}, WIREs Comput.~Mol.~Sci.~\textbf{8}, e1340 (2018).
\end{itemize}

\section{License}

\pycmf{} is distributed under the Apache License, Version~2.0 (files
\code{LICENSE} and \code{NOTICE}). It is developed by QuantumLab HCMUS,
University of Science, Vietnam National University Ho Chi Minh City.
